% Generated draft: requires mathematical and source review.
\documentclass[11pt]{article}
\usepackage[margin=25mm]{geometry}
\usepackage{fontspec}
\setmainfont{Latin Modern Roman}
\usepackage{amsmath,amssymb,mathtools,unicode-math}
\setmathfont{Latin Modern Math}
\usepackage{enumitem,needspace,xurl,hyperref}
\setlength{\parindent}{0pt}
\setlength{\parskip}{5pt}
\setlength{\emergencystretch}{3em}
\newcommand{\N}{\mathbb{N}}
\newcommand{\Z}{\mathbb{Z}}
\newcommand{\Q}{\mathbb{Q}}
\newcommand{\R}{\mathbb{R}}
\newcommand{\C}{\mathbb{C}}
\newcommand{\F}{\mathbb{F}}
\newcommand{\A}{\mathbb{A}}
\newcommand{\PP}{\mathbb{P}}
\newcommand{\E}{\mathbb{E}}
\newcommand{\eps}{\varepsilon}
\newcommand{\dd}{\,\mathrm{d}}
\begin{document}
\section*{UnsolvedMath ID: 11000197}
\textbf{Source problem code:} catalogue entry
\subsection{Definitions and mathematical statement}
Let $\pi=\pi_1(\Sigma)$ and let
\[
 X(\Sigma)=\operatorname{Hom}(\pi,\mathrm{SL}_2(\C))//\mathrm{SL}_2(\C)
\]
be the character space, with $[\rho]$ the character of a representation. The mapping class action is by precomposition, well-defined on characters. Let $U(\Sigma)$ be the image of $\operatorname{Hom}(\pi,\mathrm{SU}(2))$ in $X(\Sigma)$, and let $Q(\Sigma)$ be the quasi-Fuchsian character locus. In the usual closed genus-at-least-two setting, the latter consists of characters whose projectivizations are discrete faithful quasi-Fuchsian representations. Punctured/boundary versions need peripheral conventions that the excerpt does not provide.

\textbf{Problem 3.4.} Find $\rho\in\operatorname{Hom}(\pi,\mathrm{SL}_2(\C))$ such that, for the same orbit closure
\[
 C_\rho=\overline{\operatorname{Mod}_{\Sigma}\cdot[\rho]},
\]
both
\[
 C_\rho\cap U(\Sigma)\ne\varnothing,
 \qquad C_\rho\cap\overline{Q(\Sigma)}\ne\varnothing
\]
hold. The unitary and quasi-Fuchsian limiting characters need not coincide. The natural reading uses the ordinary complex-analytic topology, but the closure topology and precise class of surfaces are unstated and must be confirmed; Zariski closure would be a different problem. The representation itself is not required by the source to be unitary, quasi-Fuchsian, or faithful.

The appended discussion describes homological representations arising from character spaces. The chapter author and Neumann [51] are reported to obtain finite-image mapping-class actions for relative $\mathrm{SL}_2(\C)$ character varieties of the one-holed torus and four-holed sphere. Atiyah--Bott [3] are cited for compact-group topology via infinite-dimensional Morse theory; the source says the action on rational cohomology of nonsingular components factors through the symplectic representation on $H^*(\Sigma)$. Biswas [9] is cited for the resulting Torelli triviality in that nonsingular cohomological setting, not a pointwise triviality of the representation space. In contrast, Cappell--Lee--Miller [17,18] are reported to detect nontrivial Torelli action on homology of the full $\mathrm{SU}(2)$ character variety for closed $\Sigma$. These distinct homology/cohomology and smooth/full-space scopes are not collapsed into one assertion.
\subsection{Short English statement}
Find one complex surface-group character whose mapping-class orbit closure meets both the unitary-character locus and the closure of the quasi-Fuchsian locus. Specify the surface and closure topology; the source's homological-action discussion is separate context.
\subsection{Sources}
\begin{itemize}
\item Local UnsolvedMath registry, permanent ID 11000197, source code catalogue entry; no proposer or proposal year supplied.
\item Farb, edited volume \emph{Problems on Mapping Class Groups and Related Topics}, chapter ``Mapping class group dynamics on surface group representations,'' Problem 3.4, PDF page 223; \url{https://www.math.uchicago.edu/~farb/papers/mcgbook.pdf}. The registry records PDF-font-structure extraction and access on 2026-07-29; its inventory distinguishes 329 labeled headings by source page.
\item Source labels: the chapter author with Neumann [51], Atiyah--Bott [3], Biswas [9], and Cappell--Lee--Miller [17,18]. Full bibliography entries and exact theorem hypotheses are not supplied.
\end{itemize}
\end{document}
